\chapter{Заклучок}

Спроведената компаративна анализа во овој труд, овозможи сеопфатен преглед на три современи пристапи за тридимензионална реконструкција: \textbf{фотограметрија}, \textbf{невронски полиња на зрачење (NeRF)} и \textbf{гаусово распрскување}. Преку унифицирана експериментална рамка применета врз шест специјално одбрани податочни множества кои претставуваат различни реални сценарија, овој труд, има за цел да ја евалуира нивната ефикасност во однос на \textbf{визуелен квалитет}, \textbf{геометриска точност} и \textbf{временска ефикасност}, користејќи квантитативни метрики и квалитативни набљудувања на реконструираните сцени.

Резултатите од експериментите покажуваат дека \textbf{фотограметријата}, иако постар и покласичен пристап на реконструкција, сеуште обезбедува висока геометриска конзистентност и текстурна прецизност, особено кај сцени што субјектот е добро осветлен. Најголемите недостатоци кај овој метод на реконструкција се јавуваат при реконструкцијата на рефлективни и површини со ниска текстура, односно површини во кои недостасуваат клучни точки кои алгоритмот може да ги искористи при реконструкција. Дополнително, \textbf{фотограметријата} во просек има долго време на процесирање, а квалитетот на реконструкцијата е директно корелиран со преклопот на сликите од податочното множество, како и големината на самото множество.

Од друга страна, \textbf{NeRF} има значителна предност во поглед на брзината на извшување и извонредна способност за реконструкција на фините волуметриски детали и одлично претставува промени во осветлувањето на сцената. Главен недостаток на овој метод на реконструкција е потребата на податоци за тренирање на моделот, односно доколку множеството на погледи е мало, \textbf{NeRF} резултира со лошо реконстуирана сцена, како што е случајот со \textbf{античката пара}. Со тоа што \textbf{NeRF} ја претставува сцената како непрекинато волуметриско поле, некои фини детали од геометријата може да изгледаат замачкани и да допринесат до слаба реконструкција.

Како последно, \textbf{гаусовото распрскување} се истакна како најбалансирано решение за реконструкција, постигнувајќи добар компромис помеѓу \textbf{визуелен реализам при реконструкција}, \textbf{пресметковна} и \textbf{временска ефикасност}. Резултатите од експериментите покажаа добра визуелна конзистентност, често подобра од алтернативниот \textbf{NeRF}, со побрзо рендерирање на рамките и достигнување на брзини до 30 рамки во секунда.

Сумирано, анализата потврдува дека ниеден метод не доминира во сите сценарија на реконструкција. Имено, \textbf{фотограметријата} останува добар избор за технички реконструкции во контролирани услови, додека \textbf{NeRF} и \textbf{гаусовото распрскување} најдобро би реконструирале поголеми и поотворени сцени од реалниот свет. Овие резултати го нагласуваат комплементарниот карактер на анализираните методи кој отвара можност идни истражувања да се насочат кон хибридна интеграција на нервонските и геометриските принципи во еден унифициран метод на реконструкција.